\documentclass[11pt]{article}
\usepackage[a4paper,margin=1in]{geometry}
\usepackage{amsmath,amsthm,amssymb}
\usepackage{mathtools}

\newtheorem{theorem}{Theorem}[section]
\newtheorem{proposition}[theorem]{Proposition}
\newtheorem{lemma}[theorem]{Lemma}
\newtheorem{corollary}[theorem]{Corollary}
\theoremstyle{definition}
\newtheorem{definition}[theorem]{Definition}
\newtheorem{remark}[theorem]{Remark}

\newcommand{\R}{\mathbb{R}}
\newcommand{\Fr}{\mathcal{A}}
\newcommand{\Om}{\Omega}
\newcommand{\eps}{\varepsilon}
\newcommand{\rstar}{\rho_{*}}
\newcommand{\zbar}{\bar z_{\rho}}
\newcommand{\er}{e_{r,\zbar}}
\newcommand{\PiE}{\Pi(E_{\rho},\zbar)}

\title{Space-Time $\Pi$ Identity and the Integrated\\
       Quadratic Boundary Oscillation\\
       (Route~$\delta$, Plan~2, Building Block \textsc{SpaceTimeIdentity})}
\author{}
\date{}

\begin{document}
\maketitle

\begin{abstract}
We prove the space-time radial-moment identity
\[
   \int_{\rstar}^{\rho_{\delta}}\PiE\,d\mu(\rho)
   =(n-1)\!\int_{\widetilde{\Om}}\!\sigma(x)\,dx,
\]
together with the integrated quadratic Fusco--Julin boundary
oscillation
\[
   \int_{\rstar}^{\rho_{\delta}}\!\!\int_{\partial^{*}E_{\rho}}|\nu_{E_{\rho}}-\er|^{2}\,d\mathcal H^{n-1}\,d\mu(\rho)
   \le C(n,R,\rstar)\,\delta_{T}(\Om),
   \quad\textup{(B$^{2}$-int)}
\]
and the integrated radial trace
\[
   \int_{\rstar}^{\rho_{\delta}}\!\!\Bigl(\int_{\partial^{*}E_{\rho}}|H_{\zbar,\rho}-1|\,d\mathcal H^{n-1}\Bigr)^{2}d\mu(\rho)
   \le C(n,R,\rstar)\,\delta_{T}(\Om),
   \quad\textup{(T$_{1}^{2}$-int)}
\]
conditional only on the H$^{1}$-in-$\rho$ centroid bound (F) of
\cite{CentroidBound}. The proof of (T$_{1}^{2}$-int) uses a
\emph{slicing-then-squaring} decomposition $T_{1}=T_{1}^{\rm rad,slic}+T_{1}^{\rm tan}$,
in which the radial piece is bounded \emph{linearly} (not at the FMP
rate $\sqrt{\delta}$) via the good-ray equality of slicewise volumes
and the bad-ray multiplicity bound of \cite{LevelSetIdentity}; the
squared volume Fraenkel $|E_{\rho}\Delta B_{\rho}|^{2}$ is controlled
at rate $\delta$ via the new pointwise slice-rearrangement inequality
$|E_{\rho}\Delta B_{\rho}|^{2}\le C\,\PiE$. The per-$\rho$ radial $L^{1}$
trace (R) is genuinely \emph{not} used.
\end{abstract}

\section{Notation and statement of the targets}\label{sec:targets}

\subsection{Standing hypotheses}

Fix $n\ge 2$, $R\ge 2$, $\rstar\in(0,1)$. Let $\Om\subset B_R$ open
with $|\Om|=\omega_n$, $\delta:=\delta_T(\Om)\le\delta_0(n,R,\rstar)$.
Let $u=u_\Om$ be the variational torsion function,
$t:[\rstar,1]\to(0,\|u\|_\infty)$ the volume-radius parametrisation,
$E_\rho:=\{u>t(\rho)\}$, $t_B(\rho):=(1-\rho^2)/(2n)$,
$\psi(\rho):=-t'(\rho)-(-t_B'(\rho))\ge 0$,
$\rho_\delta:=1-\kappa\sqrt\delta$,
$\mu(d\rho):=(-t'(\rho))d\rho$.

For each $\rho$, $\zbar=\bar z_\rho$ is the centroid of $E_\rho$
(\cite{CentroidBound} Theorem F). By \cite{LevelSetIdentity} Lemma 2.2,
$B(\zbar,\rho/4)\subset E_\rho\subset B(\zbar,3\rho)$, in particular
$\rstar/2\le|x-\zbar|\le 2R$ on $\partial^*E_\rho$.

The L² divergence identity of \cite{LevelSetIdentity} (5.3):
\begin{equation}\label{eq:53prime}
   \int_{\partial^*E_\rho}|\nu_{E_\rho}-\er|^2 d\mathcal H^{n-1}
   = 2(P(E_\rho)-P(B_\rho))+2\PiE,
\end{equation}
where
\begin{equation}\label{eq:Pi}
   \PiE := (n-1)\int_{B_\rho(\zbar)\setminus E_\rho}|x-\zbar|^{-1}dx
         - (n-1)\int_{E_\rho\setminus B_\rho(\zbar)}|x-\zbar|^{-1}dx \ge 0.
\end{equation}

\subsection{The three target theorems}

\begin{theorem}[Space-time $\Pi$ identity, (G1)]\label{thm:G1}
$\int_{\rstar}^{\rho_\delta}\PiE d\mu(\rho)=(n-1)\int_{\widetilde\Om}\sigma(x)dx$
for a signed Fubini density $\sigma$ with
$|\sigma(x)|\le(4/n\rstar)\mathcal L^1(J_x)$, where
$J_x:=\{\rho:x\in B_\rho(\zbar)\Delta E_\rho\}$.
\end{theorem}

\begin{theorem}[Integrated (B²-int), (G3)]\label{thm:G3}
Conditional on (F),
\[
   \int_{\rstar}^{\rho_\delta}\int_{\partial^*E_\rho}|\nu_{E_\rho}-\er|^2 d\mathcal H^{n-1} d\mu(\rho)
   \le C_{B^2}(n,R,\rstar)\,\delta_T(\Om).
\]
\end{theorem}

\begin{theorem}[Integrated radial trace, (T$_1^2$-int)]\label{thm:T1sq}
Conditional on (F),
\[
   \int_{\rstar}^{\rho_\delta}\Bigl(\int_{\partial^*E_\rho}|H_{\zbar,\rho}-1|d\mathcal H^{n-1}\Bigr)^2 d\mu(\rho)
   \le C_{\rm radial}(n,R,\rstar)\,\delta_T(\Om).
\]
\end{theorem}

\section{Space-time identity (G1)}\label{sec:G1}

By \eqref{eq:Pi}, $\PiE = (n-1)\int_{\R^n}(\mathbf 1_{B_\rho(\zbar)}-\mathbf 1_{E_\rho})|x-\zbar|^{-1}dx$.
By joint Borel measurability and Fubini,
\begin{equation}\label{eq:fubini}
   \int_{\rstar}^{\rho_\delta}\PiE d\mu(\rho)
   = (n-1)\int_{\R^n}\sigma(x)dx,
\end{equation}
where $\sigma(x):=\int_{\rstar}^{\rho_\delta}(\mathbf 1_{B_\rho(\zbar)}-\mathbf 1_{E_\rho})/|x-\zbar|\,d\mu(\rho)$.
By CL containment, the integrand vanishes outside
$\widetilde\Om:=\Om\cup B_{3R}$ and $|x-\zbar|\ge\rstar/2$ on its support;
combined with $-t'\le 1/n$,
$|\sigma(x)|\le (4/n\rstar)\mathcal L^1(J_x)$. \hfill\qed

\section{Profile-gap conversion with centre drift (G2)}\label{sec:G2}

By Talenti / Nicola--Tilli profile gap \cite{LevelSetIdentity} Lemma 8.1,
$0\le t_B(\rho)-t(\rho)\le 2\delta/(\omega_n\rstar^n)$. With
$\tau_0:=\rstar/(2n)$ and on the good profile set $\{-t'\ge\tau_0\}$,
$|\rho_E(x)-\rho_B^*(x)|\le C_T\delta$ by MVT.

For moving centre, use $\varphi_x(\rho):=|x-\zbar|-\rho$;
on $\{|\bar z_\rho'|\le 1/2\}$, $\varphi_x'\le-1/2$. The good/bad
decomposition via Chebyshev on $\{|\bar z_\rho'|>1/2\}$ controlled by
(F) gives:
\begin{equation}\label{eq:G2}
   \mathcal L^1(J_x)\le C_T\delta+\int|\bar z_\rho'|\mathbf 1_{A_x}d\rho+4\int|\bar z_\rho'|^2 d\rho.
\end{equation}

\section{Assembly of (B²-int)}\label{sec:G3}

By \eqref{eq:53prime} integrated,
\[
   \int|\nu-\er|^2 dH d\mu = 2\int(P-P_\rho)d\mu+2\int\PiE d\mu.
\]
The first piece is $\le C\delta$ via Agent 1's deficit identity.
The second piece is bounded via (G1)+(G2)+(F)+Cauchy--Schwarz:
\[
   \int\PiE d\mu \le \tfrac{4(n-1)}{n\rstar}\int_{\widetilde\Om}\mathcal L^1(J_x)dx
   \le C_2(n,R,\rstar,C_F)\,\delta.
\]
Both pieces $O(\delta)$ gives (B²-int). \hfill\qed

\section{The integrated radial trace (T$_1^2$-int) --- slicing-then-squaring}\label{sec:T1}

Write $T_1=T_1(E_\rho,\zbar):=\int_{\partial^*E_\rho}||x-\zbar|/\rho-1|d\mathcal H^{n-1}$
and decompose
\[
   T_1 = T_1^{\rm rad,slic}+T_1^{\rm tan},
\]
with
\begin{align*}
   T_1^{\rm rad,slic}&:=\int_{\partial^*E_\rho}||x-\zbar|/\rho-1|\cdot|\er\cdot\nu_{E_\rho}|d\mathcal H^{n-1},\\
   T_1^{\rm tan}&:=\int_{\partial^*E_\rho}||x-\zbar|/\rho-1|\cdot(1-|\er\cdot\nu_{E_\rho}|)d\mathcal H^{n-1}.
\end{align*}

By the slicing identity \cite[Thm 18.11]{Maggi2012} applied with
$g(r):=|r/\rho-1|$:
\begin{equation}\label{eq:S1}
   T_1^{\rm rad,slic} = \int_{S^{n-1}}\sum_{j=1}^{N(\theta)}|r_{\theta,j}/\rho-1|\,r_{\theta,j}^{n-1}d\theta.
\end{equation}

\subsection{Pointwise bound for $T_1^{\rm rad,slic}$}

\begin{lemma}\label{lem:T1rad}
$T_1^{\rm rad,slic}\le C_4(n,R,\rstar)[|E_\rho\Delta B_\rho(\zbar)|+(P(E_\rho)-P(B_\rho))]$.
\end{lemma}

\begin{proof}
\emph{Good rays.} On a good ray ($N=1$), the slicewise 1D
symmetric-difference \emph{equals}:
\begin{equation}\label{eq:eq-slicewise}
   \int_0^\infty|\mathbf 1_{E_{\rho,\theta}}-\mathbf 1_{(0,\rho)}|r^{n-1}dr
   = \frac{1}{n}|r_{\theta,1}^n-\rho^n|.
\end{equation}
By MVT, $|r_{\theta,1}/\rho-1|\cdot r_{\theta,1}^{n-1}\le\beta_n|r_{\theta,1}^n-\rho^n|$
with $\beta_n=\beta_n(R,\rstar)$. Integrating and using spherical Fubini:
\[
   \int_{\{N=1\}}|r_{\theta,1}/\rho-1|r_{\theta,1}^{n-1}d\theta \le n\beta_n|E_\rho\Delta B_\rho(\zbar)|.
\]

\emph{Bad rays.} $|r_{\theta,j}/\rho-1|r_{\theta,j}^{n-1}\le (2R)^{n-1}$,
summed gives $\le (2R)^{n-1}(1+k(\theta))$. By the elementary bad-ray multiplicity bound
$\int k(\theta)\,d\theta\le(2/\rstar)^{n-1}(P(E_\rho)-P(B_\rho))$,
which follows from the slicing identity~\eqref{eq:S1}, the inner-ball
containment, and the isoperimetric inequality
(each extra crossing contributes $\ge(\rstar/2)^{n-1}$ to
$P(E_\rho)$ beyond the spherical baseline $P(B_\rho)$),

Combining gives the claim.
\end{proof}

\subsection{Pointwise bound for $T_1^{\rm tan}$}

\begin{lemma}\label{lem:T1tan}
$T_1^{\rm tan}\le (P(E_\rho)-P(B_\rho))+\PiE$.
\end{lemma}

\begin{proof}
By CL, $||x-\zbar|/\rho-1|\le 1$, so
$T_1^{\rm tan}\le\int(1-|\er\cdot\nu|)dH\le\int(1-\er\cdot\nu)dH=(P-P_\rho)+\PiE$
by \cite{LevelSetIdentity} (5.3).
\end{proof}

\subsection{Squaring and integrating}

By Lemmas \ref{lem:T1rad}, \ref{lem:T1tan} and $(a+b+c)^2\le 3(a^2+b^2+c^2)$:
\[
   T_1^2 \le 3C_5^2[|E_\rho\Delta B_\rho|^2+(P-P_\rho)^2+\PiE^2].
\]
Integrating against $d\mu$ and bounding each piece by $C\delta$:
\begin{itemize}
\item $\int|E_\rho\Delta B_\rho|^2 d\mu\le C_8\int\PiE d\mu\le C\delta$
via the slice-rearrangement (§\ref{sec:slice}) + (G3).
\item $\int(P-P_\rho)^2 d\mu\le 2P(B_\rho)\int(P-P_\rho)d\mu\le C\delta$
via Agent 1.
\item $\int\PiE^2 d\mu\le\Pi_{\max}\int\PiE d\mu\le C\delta$ since
$\Pi$ is bounded by CL.
\end{itemize}
Hence $\int T_1^2 d\mu\le C_{T_1^2}\delta$.

Combining with (G3) gives (T$_1^2$-int) since
$|H-1|\le||x-\zbar|/\rho-1|+|\er\cdot\nu-1|$ and the second term
is controlled by $|\er-\nu|$ + Cauchy--Schwarz + (G3). \hfill\qed

\section{The slice-rearrangement inequality}\label{sec:slice}

\begin{lemma}\label{lem:slice-rearrange}
$|E_\rho\Delta B_\rho(\zbar)|^2\le C_8(n,R,\rstar)\PiE$ pointwise in
the smallness regime.
\end{lemma}

\begin{proof}
\emph{Step 1: $\PiE\ge c M_1$.} Let $M_1:=\int_{E_\rho\Delta B_\rho}|\rho-|x-\zbar||dx$.
By rewriting \eqref{eq:Pi},
\[
   \PiE = (n-1)\int_{E_\rho\Delta B_\rho}\frac{|\rho-|x-\zbar||}{\rho|x-\zbar|}dx \ge\frac{n-1}{2R\rho}M_1.
\]

\emph{Step 2: $|E_\rho\Delta B_\rho|^2\le C\,M_1$.} By polar coordinates
and a slicewise Cauchy--Schwarz on $[\rstar/2,2R]$,
$\sigma(\theta)^2\le 4R\,m_1(\theta)$ with
$\sigma(\theta)=\mathcal L^1(E_{\rho,\theta}\Delta(0,\rho))$,
$m_1(\theta)=\int|r-\rho|\mathbf 1_{E_{\rho,\theta}\Delta(0,\rho)}dr$.
Cauchy--Schwarz on $S^{n-1}$ then gives
$(\int\sigma d\theta)^2\le 4R|S^{n-1}|\int m_1 d\theta$. Boosting
slice integrals to global integrals via the $r^{n-1}$ weight bounds:
$|E_\rho\Delta B_\rho|^2\le C_9\,M_1$.

Combining Steps 1--2 gives the claim.
\end{proof}

\section{Discussion: why $T_1$-route closes where $T_2$-route stalls}\label{sec:discussion}

\begin{remark}
The $T_2$-route would extract on good rays
$(r_{\theta,1}/\rho-1)^2 r_{\theta,1}^{n-1}\le C|r_{\theta,1}^{n-1}-\rho^{n-1}|$
(a square on LHS, unsigned mismatch on RHS), then need
$\int|r_{\theta,1}^{n-1}-\rho^{n-1}|d\theta\le C\sqrt{\delta_\rho}$
(FMP rate). Integrating against $d\mu$ leaks at $\sqrt\delta$.

The $T_1$-route uses linear good-ray equality \eqref{eq:eq-slicewise}
and the slice-rearrangement to get $|E_\rho\Delta B_\rho|^2\le C\Pi$,
which integrates to $C\delta$ via (G3). Squaring \emph{before}
integrating avoids the FMP bottleneck.
\end{remark}

\section{Conditionality}\label{sec:conditionality}

\textbf{Inputs used (all unconditional given (F)):}
CL containment; slicing identity (Maggi); bad-ray multiplicity bound (elementary, see proof of Lemma~\ref{lem:T1rad}); slicewise good-ray equality; spherical Fubini;
divergence identity (5.3); Talenti profile gap;
slice-rearrangement Lemma \ref{lem:slice-rearrange}; (F) of
\cite{CentroidBound}.

\textbf{Inputs explicitly not used:} the per-$\rho$ pointwise radial-trace inputs (1.1), (1.2)=(R),
(2.4); pointwise (B²); Agent 3 (7.1); any selection, graph entry,
Schauder, Fuglede.

\begin{thebibliography}{99}

\bibitem{LevelSetIdentity}
Plan~2 building block \textsc{LevelSetIdentity}.

\bibitem{MetricFramework}
Plan~2 building block \textsc{MetricFramework}.

\bibitem{CentroidBound}
Plan~2 building block \textsc{CentroidBound}.

\bibitem{WeightedMetricTrace}
Plan~2 building block \textsc{WeightedMetricTrace}.

\bibitem{Maggi2012}
F.~Maggi, Cambridge Univ. Press, 2012. Thm 15.9, 18.11, 19.4.

\bibitem{AFP2000}
L.~Ambrosio, N.~Fusco, D.~Pallara, Oxford, 2000.

\bibitem{FMP2010}
A.~Figalli, F.~Maggi, A.~Pratelli, Invent. Math. \textbf{182} (2010), 167--211.

\bibitem{FJ2014}
N.~Fusco, V.~Julin, Calc. Var. PDE \textbf{50} (2014), 925--937.

\bibitem{CL2012}
M.~Cicalese, G.\,P.~Leonardi, ARMA \textbf{206} (2012), 617--643.

\bibitem{AFM2013}
E.~Acerbi, N.~Fusco, M.~Morini, CMP \textbf{322} (2013), 515--557.

\bibitem{Talenti1976}
G.~Talenti, Ann. Mat. Pura Appl. (4) \textbf{110} (1976), 353--372.

\bibitem{NicolaTilli}
F.~Nicola, P.~Tilli, Invent. Math. \textbf{230} (2022), 1--30.

\end{thebibliography}

\end{document}
